\chapter{Desarrollo del Trabajo}

En esta fase se ha definido el problema como una búsqueda de un candidato de subclave con puntuación diferencial elevada. La puntuación cuenta cuántos de los pares disponibles superan un filtro después de aplicar el descifrado parcial asociado a un candidato. Esta formulación conecta el procedimiento clásico de conteo con la búsqueda cuántica de extremos descrita en la literatura \citep{zhou2019}.

\section{Flujo del procedimiento}

El flujo previsto consta de las siguientes operaciones:
\begin{enumerate}
 \item preparar el conjunto de pares con la diferencia de entrada requerida;
 \item definir para cada candidato el descifrado parcial de la última ronda;
 \item comprobar el filtro diferencial y agregar los votos para formar \(A(k)\);
 \item comparar candidatos y seleccionar uno con puntuación máxima;
 \item verificar el candidato recuperado mediante los pares y el modelo de cifrado.
\end{enumerate}

En el método clásico se calcula el conteo para cada candidato. En la variante cuántica, la puntuación debe calcularse de forma coherente y reversible para ser utilizada en las comparaciones iterativas de la búsqueda de máximo. El diseño lógico debe especificar los registros de entrada, salida y trabajo, así como el procedimiento para descomputar la información temporal.

\section{Estado de implementación}

El repositorio indicado por los autores se considera el artefacto que deberá permitir reproducir la parte clásica y validar la función de puntuación. En esta versión del borrador aún no se ha documentado una inspección de los ficheros ni una ejecución reproducible del código; por ello, no se atribuye al repositorio ningún resultado numérico concreto. Como parte del desarrollo incremental del TFE, se completará una tabla con lenguaje y dependencias, comando de ejecución, parámetros, salidas esperadas y pruebas con casos pequeños cuya puntuación pueda comprobarse manualmente.

\section{Plan de evaluación}

La evaluación medirá (i) la corrección del conteo respecto de una implementación de referencia; (ii) la tasa de recuperación del candidato correcto en instancias con clave conocida; (iii) el número de pares \(N\), el tamaño de candidato \(b\) y el tiempo de ejecución clásico; y (iv) la estimación de consultas cuánticas y coste de una llamada reversible. Se separarán las mediciones observadas de las cotas asintóticas y se indicarán límites de extrapolación.
